\chapter{Clearing and Resolution}
\label{ch:clearing}

\section*{The Clearing Mode}

\begin{versebox}
When the question has been answered ---\\
truly answered, not just glossed ---\\
the path begins in earnest;\\
until then, the point is lost.
\end{versebox}

\noindent Have you ever asked someone a question and gotten an answer that \textit{sounded} like an answer but actually just restated the problem in different words? You ask your doctor, ``Why do I feel tired all the time?'' and she says, ``You have chronic fatigue.'' You ask a mechanic, ``Why won't my car start?'' and he says, ``The engine won't turn over.'' In each case, you got a response --- but not a resolution.\footnote{Pali: \textit{``Tattha katamo sodhano hāro? `Vissajjitamhi pañhe'\,''tigāthā.''} The mnemonic fragment is ``when the question has been resolved'' (\textit{vissajjitamhi pañhe}). The word \textit{sodhana} means ``clearing, purifying, cleansing'' --- the mode that distinguishes a genuine resolution from a mere restatement. A question is ``cleared'' (\textit{suddha}) only when the answer provides genuinely new information --- what the Netti calls a ``clean beginning'' (\textit{suddho ārambho}).}

The Clearing Mode is the tool that tells the difference. It distinguishes between answers that merely \textit{clear the term} --- substitute one word for another --- and answers that actually \textit{resolve} the question by providing a genuine ``clean beginning'': new information that moves you forward.

\section*{Ajita's Questions}

The demonstration uses one of the most famous exchanges in the Canon --- the dialogue between the Venerable Ajita and the Buddha from the Pārāyana, the ``Way to the Far Shore.''\footnote{Pali: \textit{``Yathā āyasmā ajito pārāyane bhagavantaṃ pañhaṃ pucchati.''} The Pārāyana (Sn 5) is one of the oldest strata of the Pali Canon, a collection of questions posed to the Buddha by sixteen Brahmin students. Ajita's questions come first (Sn 1032--1039).}

Ajita asks four questions in a single verse:

\begin{versebox}
By what is the world enveloped?\\
By what does it not shine?\\
What do you call its smearing?\\
What is its great fear?
\end{versebox}

And the Buddha answers:

\begin{versebox}
By ignorance the world is enveloped, Ajita ---\\
by desire and heedlessness it does not shine.\\
I call yearning its smearing;\\
suffering is its great fear.
\end{versebox}

Four questions, four answers. But the Clearing Mode notices something crucial. The first three answers merely \textit{clear the term}. ``By what is the world enveloped?'' --- ``By ignorance.'' That replaces one unknown (``by what'') with another (``ignorance''), but it does not resolve the question. You still need to know what ignorance is, how it works, what to do about it. The same with desire and yearning --- they are identifications, not resolutions.\footnote{Pali: \textit{``Kenassu nivuto loko\,''ti pañhe `avijjāya nivuto loko\,''ti bhagavā padaṃ sodheti, no ca ārambhaṃ. `Kenassu nappakāsatī\,''ti pañhe `vivicchā pamādā nappakāsatī\,''ti bhagavā padaṃ sodheti, no ca ārambhaṃ. `Kissābhilepanaṃ brūsī\,''ti pañhe `jappābhilepanaṃ brūmī\,''ti bhagavā padaṃ sodheti, no ca ārambhaṃ.''} Three term-clearings (\textit{padaṃ sodheti}) without a clean beginning (\textit{no ca ārambhaṃ}). ``By what is the world enveloped?'' $\to$ ``By ignorance'' = term cleared, but no clean beginning. ``By what does it not shine?'' $\to$ ``By desire and heedlessness'' = term cleared, no clean beginning. ``What do you call its smearing?'' $\to$ ``I call yearning its smearing'' = term cleared, no clean beginning.}

But the fourth answer is different. ``What is its great fear?'' --- ``Suffering is its great fear.'' \textit{That} is a clean beginning. Why? Because ``suffering'' is not just another synonym. It is the First Noble Truth --- the starting point of the entire path. When the Buddha says ``suffering,'' he opens a door that leads somewhere: to the origin, the cessation, the path. The first three answers replaced labels. The fourth launched a journey.\footnote{Pali: \textit{``Kiṃsu tassa mahabbhaya\,''nti pañhe `dukkhamassa mahabbhaya\,''nti suddho ārambho.''} ``Suffering is its great fear'' = clean beginning (\textit{suddho ārambho}). The Netti's distinction is precise: a ``term-clearing'' substitutes an equivalent word; a ``clean beginning'' provides the anchor-point from which the teaching can unfold. ``Suffering'' is a clean beginning because it is not just a label --- it is the First Noble Truth, which necessarily implies the other three truths and the entire practice leading to liberation.}

\section*{A Second Round}

Ajita presses on:

\begin{versebox}
Currents flow everywhere ---\\
what is their obstruction?\\
Tell me the restraint of currents:\\
by what are currents dammed?
\end{versebox}

The Buddha:

\begin{versebox}
Whatever currents are in the world ---\\
mindfulness is their obstruction.\\
I call mindfulness the restraint of currents;\\
by wisdom they are dammed.
\end{versebox}

Again, the first two answers clear the term: ``What is their obstruction?'' --- ``Mindfulness.'' ``What is their restraint?'' --- ``Mindfulness.'' Substitutions. But the last answer --- ``by wisdom they are dammed'' --- is the clean beginning. Wisdom is not just another word for restraint; it is the faculty that actually stops the flow. It resolves the question.\footnote{Pali: \textit{``Savanti sabbadhi sotā, sotānaṃ kiṃ nivāraṇa\,''nti pañhe `yāni sotāni lokasmiṃ, sati tesaṃ nivāraṇa\,''nti bhagavā padaṃ sodheti, no ca ārambhaṃ. `Sotānaṃ saṃvaraṃ brūhi, kena sotā pidhīyare\,''ti pañhe `sotānaṃ saṃvaraṃ brūmi, paññāyete pidhīyare\,''ti suddho ārambho.''} First two answers: term cleared, no clean beginning. Final answer: ``by wisdom they are dammed'' = clean beginning.}

\section*{The Final Question}

And then the deepest question of all:

\begin{versebox}
Wisdom and mindfulness,\\
and name-and-form, sir ---\\
tell me, when asked:\\
where do these cease?
\end{versebox}

The Buddha:

\begin{versebox}
This question you have asked, Ajita ---\\
I will answer you.\\
Where name-and-form ceases\\
without remainder ---\\
by the cessation of consciousness,\\
there they cease.
\end{versebox}

This entire answer is a clean beginning. No term-clearing at all. The Buddha goes straight to the resolution: consciousness ceases, and with it, name-and-form ceases. That is where everything stops. That is the answer.\footnote{Pali: \textit{``Yametaṃ pañhaṃ apucchi, ajita taṃ vadāmi te; / Yattha nāmañca rūpañca, asesaṃ uparujjhati; / Viññāṇassa nirodhena, etthetaṃ uparujjhatī\,''ti. Suddho ārambho.''} The entire answer is a clean beginning. No preliminary term-clearing. The Buddha answers directly: name-and-form ceases with the cessation of consciousness (\textit{viññāṇassa nirodhena}).}

The principle of the Clearing Mode: where there is a clean beginning, the question has been resolved. Where there is no clean beginning --- where the answer merely swaps one term for another --- the question remains open, the work unfinished.\footnote{Pali: \textit{``Yattha evaṃ suddho ārambho, so pañho visajjito bhavati. Yattha pana ārambho asuddho, na tāva so pañho visajjito bhavati. Tenāha āyasmā mahākaccāyano `vissajjitamhi pañhe\,''ti.''} The Netti's closing formula for the Clearing Mode: ``Where there is a clean beginning, the question is resolved. Where the beginning is not clean, the question is not yet resolved.'' This is a hermeneutical razor: it separates genuine answers from sophisticated-sounding non-answers.}

This is a razor for reading. When you encounter a sutta where the Buddha seems to answer a question, ask: did he actually resolve it, or did he just clear a term? If the latter, the real answer is still coming.

\ornament

\section*{The Resolution Mode}

\begin{versebox}
States of one nature and many ---\\
the general and the specific.\\
Match your answer to the question:\\
every level is forensic.
\end{versebox}

\noindent Imagine you are looking at a map. At one zoom level, you see ``Europe.'' Zoom in, and you see France, Germany, Italy. Zoom in further, and you see Paris, Berlin, Rome. Zoom in further still, and you see streets, buildings, rooms.\footnote{Pali: \textit{``Tattha katamo adhiṭṭhāno hāro? `Ekattatāya dhammā, yepi ca vemattatāya niddiṭṭhā\,''ti.''} The mnemonic fragment is ``states that are of one nature, and those designated as of different nature'' (\textit{ekattatāya dhammā, yepi ca vemattatāya niddiṭṭhā}). The word \textit{adhiṭṭhāna} means ``resolution, determination, fixing'' --- the mode that ``resolves'' a general term into its specific components, or recognizes that a cluster of specifics belongs to a single general category. The two key terms are \textit{ekattatā} (``oneness, sameness of nature'' --- the general category) and \textit{vemattatā} (``difference, diversity'' --- the specific subcategories).}

The Resolution Mode is about zoom level. Every term the Buddha used has a general meaning --- its ``oneness'' --- and a specific meaning --- its ``difference.'' Knowing which level a teaching is operating at, and matching your understanding to that level, is the Resolution Mode.

\section*{Suffering: One and Many}

``Suffering'' is one thing. That is the general level --- oneness. But ask ``What is suffering?'' and the answer fans out into many: birth is suffering, aging is suffering, sickness is suffering, death is suffering, being united with what you dislike is suffering, being separated from what you love is suffering, not getting what you want is suffering --- and in brief, the five aggregates of clinging are suffering.\footnote{Pali: \textit{``Dukkha\,''nti ekattatā. Tattha katamaṃ dukkhaṃ? Jāti dukkhā, jarā dukkhā, byādhi dukkho, maraṇaṃ dukkhaṃ, appiyehi sampayogo dukkho, piyehi vippayogo dukkho, yampicchaṃ na labhati tampi dukkhaṃ, saṃkhittena pañcupādānakkhandhā dukkhā, rūpā dukkhā, vedanā dukkhā, saññā dukkhā, saṅkhārā dukkhā, viññāṇaṃ dukkhaṃ. Ayaṃ vemattatā.''} ``Suffering'' = oneness (\textit{ekattatā}). Its enumeration into birth, aging, disease, death, the five types of painful conjunction/disjunction, and the five aggregates of clinging = difference (\textit{vemattatā}). The general term contains all the specifics; the specifics elaborate the general term.}

One word. Twelve specifics. Same truth, different zoom levels.

The pattern repeats for every major term. ``The origin of suffering'' is one thing at the general level. But zoom in, and it is the craving that makes for renewed existence --- specifically, sensual craving, craving for existence, and craving for non-existence.\footnote{Pali: \textit{``Dukkhasamudayo\,''ti ekattatā. Tattha katamo dukkhasamudayo? Yāyaṃ taṇhā ponobhavikā nandīrāgasahagatā tatratatrābhinandinī. Seyyathidaṃ, kāmataṇhā bhavataṇhā vibhavataṇhā. Ayaṃ vemattatā.''}} ``Cessation'' is one thing --- but zoom in, and you find reflective cessation, non-reflective cessation, cessation of attraction, cessation of aversion, cessation of conceit, cessation of spite, cessation of envy, cessation of stinginess, cessation of all defilements.\footnote{Pali: \textit{``Nirodho\,''ti ekattatā. Tattha katamo nirodho? Paṭisaṅkhānirodho appaṭisaṅkhānirodho anunayanirodho paṭighanirodho mānanirodho makkhanirodho paḷāsanirodho issānirodho macchariyanirodho sabbakilesanirodho. Ayaṃ vemattatā.''} Two major types of cessation: reflective cessation (\textit{paṭisaṅkhānirodha}) --- achieved through deliberate contemplation --- and non-reflective cessation (\textit{appaṭisaṅkhānirodha}) --- the natural fading of defilements. Then eight more specific cessations, each naming a particular defilement that ceases.} ``The path'' is one thing --- but zoom in, and there are paths to hell, to the animal realm, to the ghost realm, to the human realm, to heaven, and to nirvana.\footnote{Pali: \textit{``Maggo\,''ti ekattatā. Tattha katamo maggo? Nirayagāmī maggo tiracchānayonigāmī maggo pettivisayagāmī maggo \ldots nibbānagāmī maggo. Ayaṃ vemattatā.''} A startling reminder: ``path'' (\textit{magga}) does not always mean the Noble Eightfold Path. There are paths leading downward too. The Resolution Mode prevents you from assuming that every mention of ``path'' in a sutta refers to the same thing.}

\section*{Forty-Two Ways to See a Body}

The Resolution Mode's most vivid demonstration is the analysis of ``form.'' At the general level, form is one thing. But zoom in, and it resolves into the four great elements --- earth, water, fire, and wind --- and whatever form is derived from them. Zoom in further, and each element resolves into its bodily manifestations.\footnote{Pali: \textit{``Rūpa\,''nti ekattatā. Tattha katamaṃ rūpaṃ? Cātumahābhūtikaṃ rūpaṃ catunnaṃ mahābhūtānaṃ upādāya rūpassa paññatti. Tattha katamāni cattāri mahābhūtāni? Pathavīdhātu āpodhātu tejodhātu vāyodhātu.''} ``Form'' = oneness. Four great elements + derived form = difference. Each element is then further resolved into its bodily aspects.}

The earth element appears in twenty aspects of the body: head-hair, body-hair, nails, teeth, skin, flesh, sinews, bones, bone-marrow, kidneys, heart, liver, diaphragm, spleen, lungs, large intestine, small intestine, stomach contents, feces, and brain. The water element appears in twelve: bile, phlegm, pus, blood, sweat, fat, tears, grease, saliva, snot, joint-fluid, and urine. The fire element in four: what warms, what ages, what burns, what digests. The wind element in six: upward-going winds, downward-going winds, belly-winds, bowel-winds, winds coursing through the limbs, and the breath.\footnote{Pali: \textit{``Vīsatiyā ākārehi pathavīdhātuṃ vitthārena pariggaṇhāti \ldots dvādasahi ākārehi āpodhātuṃ \ldots catūhi ākārehi tejodhātuṃ \ldots chahi ākārehi vāyodhātuṃ \ldots''} The full forty-two-part analysis of the body: earth (20) + water (12) + fire (4) + wind (6) = 42 aspects. This analysis comes from the Mahāhatthipadopama Sutta (MN 28) and the Mahārāhulovāda Sutta (MN 62). It is one of the standard meditation practices for seeing the body as a collection of elements rather than a unified self.}

Forty-two aspects. And when you examine them one by one --- weighing, investigating, scrutinizing --- you find nothing worth grasping. The Netti compares it to searching through a muddy ditch, a rubbish heap, a latrine, a charnel ground. You sift through the whole thing, and there is nothing there to hold onto.\footnote{Pali: \textit{``Evaṃ imehi dvācattālīsāya ākārehi vitthārena dhātuyo sabhāvato upalakkhayanto tulayanto parivīmaṃsanto pariyogāhanto paccavekkhanto na kiñci gayhūpagaṃ passati kāyaṃ vā kāyapadesaṃ vā, yathā candanikaṃ pavicinanto na kiñci gayhūpagaṃ passeyya, yathā saṅkāraṭṭhānaṃ pavicinanto \ldots yathā vaccakuṭiṃ pavicinanto \ldots yathā sivathikaṃ pavicinanto na kiñci gayhūpagaṃ passeyya.''} Four similes for the fruitless search through the body: (1) searching through a muddy ditch (\textit{candanikā}), (2) a rubbish heap (\textit{saṅkāraṭṭhāna}), (3) a latrine (\textit{vaccakuṭi}), (4) a charnel ground (\textit{sivathika}). In each case: you search and find nothing worth taking. The body, analyzed into its forty-two component elements, yields nothing that could serve as a self.}

And the conclusion? ``This is not mine. This is not what I am. This is not my self.'' Seeing it with right wisdom, you become disenchanted. You free the mind from the earth element, the water element, the fire element, the wind element. One word --- ``form'' --- resolved into forty-two parts, each one empty of self.\footnote{Pali: \textit{``Taṃ `netaṃ mama, nesohamasmi, na meso attā\,''ti evametaṃ yathābhūtaṃ sammappaññāya daṭṭhabbaṃ, evametaṃ yathābhūtaṃ sammappaññāya disvā pathavīdhātuyā nibbindati, pathavīdhātuyā cittaṃ virājeti \ldots vāyodhātuyā nibbindati, vāyodhātuyā cittaṃ virājeti. Ayaṃ vemattatā.''} From MN 28. The practitioner applies the three characteristics (``not mine, not I, not my self'') to each element in turn. Disenchantment (\textit{nibbidā}) and dispassion (\textit{virāga}) follow naturally. This is the Resolution Mode applied as meditation practice: the general term ``form'' is resolved into specifics, and each specific is seen as not-self.}

\section*{The Many Faces of Ignorance --- and Wisdom}

The Resolution Mode applies the same zoom-lens to the mental side. ``Ignorance'' is one thing at the general level. But zoom in, and it is not-knowing suffering, not-knowing its origin, not-knowing its cessation, not-knowing the path. Not-knowing the past, the future, the past-and-future. Not-knowing dependent origination. And then a cascade of synonyms: not-seeing, not-understanding, not-awakening, not-penetrating, not-discerning, not-examining, not-being-present-to, foolishness, stupidity, lack of clear comprehension, delusion, confusion, bewilderment, ignorance-as-flood, ignorance-as-bond, ignorance-as-underlying-tendency, ignorance-as-obsession, ignorance-as-barrier --- and delusion itself, the unwholesome root.\footnote{Pali: \textit{``Avijjā\,''ti ekattatā. Tattha katamā avijjā? Dukkhe aññāṇaṃ \ldots idappaccayatāpaṭiccasamuppannesu dhammesu aññāṇaṃ, yaṃ evarūpaṃ aññāṇaṃ adassanaṃ anabhisamayo ananubodho asambodho appaṭivedho asallakkhaṇā anupalakkhaṇā apaccupalakkhaṇā asamavekkhaṇaṃ apaccakkhakammaṃ dummejjhaṃ bālyaṃ asampajaññaṃ moho pamoho sammoho avijjā avijjogho avijjāyogo avijjānusayo avijjāpariyuṭṭhānaṃ avijjālaṅgī moho akusalamūlaṃ. Ayaṃ vemattatā.''} ``Ignorance'' = oneness. Its resolution into eight types of not-knowing (the Four Truths, past, future, past-and-future, dependent origination) plus at least twenty synonyms = difference. The sheer number of synonyms underlines how pervasive ignorance is --- it infiltrates every aspect of the mind.}

And ``true knowledge''? One thing at the general level --- but zoom in, and it fans out into the same eight types of knowing, plus an even longer list: wisdom, discernment, investigation, discrimination, examination, learning, skill, subtlety, the power of analysis, reflection, breadth, intelligence, guidance, insight, clear comprehension, the goad, the wisdom-faculty, the wisdom-power, the wisdom-sword, the wisdom-tower, the wisdom-light, the wisdom-radiance, the wisdom-torch, the wisdom-jewel, non-delusion, investigation-of-states, right view, the awakening-factor of investigation-of-states, a path-factor, included in the path.\footnote{Pali: \textit{``Vijjā\,''ti ekattatā. Tattha katamā vijjā? Dukkhe ñāṇaṃ \ldots yā evarūpā paññā pajānanā vicayo pavicayo dhammavicayo saṃlakkhaṇā upalakkhaṇā paccupalakkhaṇā paṇḍiccaṃ kosallaṃ nepuññaṃ vebhabyā cintā upaparikkhā bhūrī medhā pariṇāyikā vipassanā sampajaññaṃ patodo paññā paññindriyaṃ paññābalaṃ paññāsatthaṃ paññāpāsādo paññāāloko paññāobhāso paññāpajjoto paññāratanaṃ amoho dhammavicayo sammādiṭṭhi dhammavicayasambojjhaṅgo maggaṅgaṃ maggapariyāpannaṃ. Ayaṃ vemattatā.''} The list of wisdom's aspects includes architectural metaphors (tower, \textit{pāsāda}), light metaphors (radiance, torch, light), martial metaphors (sword, \textit{sattha}), and economic metaphors (jewel, \textit{ratana}). Each names a different function or manifestation of the same underlying capacity.}

The zoom-out reveals one thing. The zoom-in reveals dozens. Neither level is wrong. But confusing them --- reading a general teaching as if it were specific, or a specific teaching as if it were general --- produces misunderstanding.

\section*{The Principle}

The Resolution Mode also resolves ``meditator'' into trainee, non-trainee, and neither-trainee-nor-non-trainee; ``concentration'' into more than twenty types; ``practice'' into painful-and-slow, painful-and-quick, pleasant-and-slow, and pleasant-and-quick; ``body'' into the physical body (thirty-two parts) and the mental body (feeling, perception, volition, mind, contact, and attention).\footnote{Pali: The Netti resolves \textit{jhāyī} (``meditator'') into eight types, \textit{samādhi} (``concentration'') into twenty-four types, \textit{paṭipadā} (``practice'') into eleven types, and \textit{kāya} (``body'') into two main divisions with numerous sub-categories. \textbf{A note on our approach}: the original Pali presents each of these as a separate \textit{ekattatā/vemattatā} (oneness/difference) pair, listing all subcategories. We have summarized the list here; the full enumerations are preserved in this note for reference. See: \textit{jhāyī} = sekkho / asekkho / nevasekkhānāsekkho / ājāniyo / assakhaluṅko / diṭṭhuttaro / taṇhuttaro / paññuttaro; \textit{samādhi} = saraṇo / araṇo / savero / avero / sabyāpajjo / abyāpajjo / sappītiko / nippītiko (and sixteen more); \textit{paṭipadā} = āgāḷha / nijjhāma / majjhima / akkhamā / khamā / samā / damā / dukkhā dandhābhiññā / dukkhā khippābhiññā / sukhā dandhābhiññā / sukhā khippābhiññā.}

The rule is simple: match the level of your answer to the level of the question. If someone asks about oneness, answer with oneness. If someone asks about difference, answer with difference. If someone asks about beings, answer about beings. If someone asks about mental qualities, answer about mental qualities. However the question is framed, that is how the answer should be framed.\footnote{Pali: \textit{``Evaṃ sutte vā veyyākaraṇe vā gāthāyaṃ vā pucchitena vīmaṃsayitabbaṃ, kiṃ ekattatāya pucchati, udāhu vemattatāyāti. Yadi ekattatāya pucchitaṃ, ekattatāya visajjayitabbaṃ. Yadi vemattatāya pucchitaṃ, vemattatāya visajjayitabbaṃ. Yadi sattādhiṭṭhānena pucchitaṃ, sattādhiṭṭhānena visajjayitabbaṃ. Yadi dhammādhiṭṭhānena pucchitaṃ, dhammādhiṭṭhānena visajjayitabbaṃ. Yathā yathā vā pana pucchitaṃ, tathā tathā visajjayitabbaṃ. Tenāha āyasmā mahākaccāyano `ekattatāya dhammā\,''ti.''} The hermeneutical rule: match the answer's frame to the question's frame. Four types of resolution: by oneness (\textit{ekattatā}), by difference (\textit{vemattatā}), by beings (\textit{sattādhiṭṭhāna}), and by phenomena (\textit{dhammādhiṭṭhāna}). ``However it is asked, thus it should be answered.'' This is a remarkably modern hermeneutical principle: interpretation must respect the register and level of the text being interpreted.}

The Clearing Mode tells you whether a question has been truly answered. The Resolution Mode tells you at what level it is being answered. Together, they prevent two of the most common errors in reading: mistaking a restatement for a resolution, and mistaking the zoom level of the teaching.
